\section{Results}\label{sec:q:res}

Every table in this section is one solved path per cell, at the planner corner of the policy space unless the dials are stated, from the calibrated set and initial stocks of Section~\ref{sec:q:data}, at the horizon $T=400$ that Section~\ref{sec:q:res:baseline} shows to be settled. Dates are periods since 1900. The metals-only bound of Section~\ref{sec:q:data:metals} is not run: at its stock-effect elasticity of zero the reserve pays no dividend, nothing in the residual system of Section~\ref{sec:q:model:residuals} holds the reserve non-negative, and the solved reserve crosses zero in 1960 and ends far below it. The solver converges on that system because the system is consistent; the economy it describes has no scarcity, and Section~\ref{sec:q:res:baseline} says what the smallest repair is.

\subsection{The calibrated baseline}\label{sec:q:res:baseline}

\smalltitle{The state.} Table~\ref{tab:q:res:taxonomy} places the baseline, the soft ceiling without a floor, in state C, perpetual circular growth, by the closure criterion of \theory{Section}{The circular balanced growth path}: the cumulated leakage sum of that criterion is $0.08$, and the effective survival factor of the waste-stock costate, the diagnostic Section~\ref{sec:q:model:residuals} names, is $0.9995$ at the horizon against the primitive $1-\mu=0.965$. The markers of the state that a finite path can show all hold. The treated share is zero until 2003, is on its calibration target in 2015, and is one from 2072; the product $a\varpi$ is $0.998$ at the horizon; recycling capital per treated tonne rises linearly at the rate $g/\xi$ the exponential tail predicts; and cumulated leakage is $0.08$ of the material budget, so the cumulated leakage lemma of \theory{Section}{The material budget and four accounting facts} holds with room.

%% GENERATED by model/scripts/run_experiments.jl on 2026-09-18
%% Do not edit by hand: the file is rewritten by the next run.
\begin{table}[!htb]
\centering
\begin{threeparttable}
\caption{Where the economy lands in the taxonomy}\label{tab:q:res:taxonomy}
\begin{tabular}{lcrrrrr}
\toprule
Case & State & $\mathcal M_\infty$ & $\mathcal T$ & $\mathcal M_\infty/(\bar R\mathcal T)$ & Gate fee $<0$ & $T$ reached \\
\midrule
abar 1, Rbar 0 & C & 86010.32 & 42.9 & $\infty$ & 139 & 400 \\
abar 1, Rbar 9.5 & C & 85387.41 & 42.9 & 209.67 & 132 & 400 \\
abar 1, Rbar 23.8 & C & 84466.88 & 42.9 & 82.95 & 117 & 400 \\
abar 1, Rbar 47.5 & - & -- & -- & -- & -- & 150 \\
abar 0.708, Rbar 0 & B & 42586.85 & 42.7 & -- & 171 & 400 \\
abar 0.708, Rbar 9.5 & A & 42207.83 & 42.7 & -- & 166 & 400 \\
abar 0.708, Rbar 23.8 & A & 41607.92 & 42.7 & -- & 157 & 400 \\
abar 0.708, Rbar 47.5 & - & -- & -- & -- & -- & 150 \\
\bottomrule
\end{tabular}
\begin{tablenotes}[flushleft]
\footnotesize
\item[] \textit{Note:} States are those of the theory note's taxonomy: A collapse, B balanced dematerialization, C perpetual circular growth. $\mathcal T$ is the residence time $1/\mu+\sigma_\infty/\delta$. The survival ratio is infinite without a floor and is reported only under a soft ceiling $\bar a = 1$; under a hard ceiling the state is the ceiling's, the loop being unable to close whatever the ratio, and the cell carries a dash. Dates are periods since the base year; a dash in a date column is an event that does not occur within the horizon reached. Where a floor meets a hard ceiling the loop cannot close and the material era has bounded duration, by the theory note's Lemma on finite cumulative throughput: cumulative handling cannot exceed the material budget over $1-\bar a$, which at $\bar R$ a period is between $2.5\times10^{4}$ and $6.2\times10^{4}$ periods in the collapse cells of this table. That bound, and not a computed date, is what those cells say about the length of the era; no shutdown date is reported, because the search's objective is monotone in the date and it returns the last date its branch solves.
\end{tablenotes}
\end{threeparttable}
\end{table}


\smalltitle{The horizon.} The horizon is settled by the first criterion of Section~\ref{sec:q:sol:horizon}: on every solved row of Table~\ref{tab:q:res:taxonomy} the 1900 to 2000 window moves by less than $3\times10^{-5}$ in output and consumption when the horizon is extended from $200$ to $400$, against the tolerance of $10^{-3}$. The residual of every row is at $10^{-10}$, the ledger holds to $10^{-12}$, the market and planner residuals agree to $10^{-14}$, the transversality factor $\beta e^{(1-\eta)g}$ is below one at the value Table~\ref{tab:q:res:sensitivity} gives, and the four long-run conditions of \theory{Section}{Scope: maintained long-run assumptions} hold at the circular rate. The conditioning wall of Section~\ref{sec:q:sol:horizon} is not reached: the reserve at the horizon is $0.18$ of its 1900 level, and it does not fall to two percent of it within the horizon on any row. The two rows at the largest floor of the grid do not solve by any route, and Section~\ref{sec:q:res:taxonomy} says why.

\smalltitle{Sufficiency.} The convexity conditions of \theory{Appendix}{Sufficiency} fail on the same four counts as on the illustrative set, which Section~\ref{sec:q:data:checks} records, so the computed path is verified rather than proved optimal. The feasible-direction test of Section~\ref{sec:q:sol:verification} at $T=300$ draws $400$ perturbations, of which $90$ are feasible, and finds a best gain of $-3\times10^{-10}$; the extraction and investment profiles are single-peaked in the perturbation; and the shutdown is absorbing for a depleted economy and not for the terminal state of the growth path, which is the ordering \theory{Section}{The operating set, and whether shutdown is absorbing} requires.

\smalltitle{The gate fee.} Along the baseline the gate fee is positive at the base year, rises to a peak in the first decade of this century, crosses zero on the date Table~\ref{tab:q:res:taxonomy} gives, and is negative and falling thereafter: the waste stock is a liability while the pollution it releases is the larger term, and an asset once the material it holds is. Section~\ref{sec:q:res:gatefee} gives the path in dollars per tonne and its range.

\smalltitle{The path against the century.} The calibration reads every parameter from a share, a rate or a single year, so the model's own path over 1900 to 2015 is a test of it rather than a fit, and Table~\ref{tab:q:res:fit} reports the test. The path does not reproduce the century's levels. Output in 2015 is three quarters of the data's, material input three fifths, extraction a little over half, and secondary input nearly twice. The miss begins in the base year, where the planner chooses a quarter less material than the data show and $A_0$ was set with: at the data's 1900 input the marginal product of material at the calibrated value share is 46 dollars per tonne while the calibrated marginal extraction cost there is about 64, so the planner uses less, and a unit elasticity of substitution cannot carry an intensity that fell by a factor of $2.4$ over a century in which the real price index of Table~\ref{tab:q:data:extraction} rose by less than a factor of two. The output miss then compounds it, because the composite trend of Section~\ref{sec:q:data:production} was read off the data's capital and material paths and the model's material path is lower throughout. This is the limitation of a calibration fitted to shares and rates and not to levels; a calibration that matched the 2015 levels would move $A_0$, $\varsigma$ and the trend jointly, and is not attempted here.

%% GENERATED by model/scripts/run_experiments.jl on 2026-09-18
%% Do not edit by hand: the file is rewritten by the next run.
\begin{table}[!htb]
\centering
\begin{threeparttable}
\caption{The calibrated path against the century it was calibrated on}\label{tab:q:res:fit}
\begin{tabular}{llrrrrrrrr}
\toprule
Series & Unit & \multicolumn{2}{c}{1900} & \multicolumn{2}{c}{1950} & \multicolumn{2}{c}{2000} & \multicolumn{2}{c}{2015} \\
 &  & Model & Data & Model & Data & Model & Data & Model & Data \\
\midrule
$Y$, output & trillion \$ & 3.4 & 3.5 & 13.1 & 8.5 & 52.9 & 60.0 & 80.3 & 107.2 \\
$K$, capital & trillion \$ & 10.6 & 10.6 & 35.0 & 30.5 & 142.9 & 208.8 & 222.1 & 448.2 \\
$R$, material input & Gt & 5.7 & 7.6 & 17.1 & 15.4 & 42.8 & 60.3 & 59.1 & 95.0 \\
$N$, extraction & Gt & 5.7 & 7.3 & 17.1 & 15.0 & 42.8 & 57.2 & 47.7 & 88.9 \\
$R^R$, secondary input & Gt & 0.0 & 0.3 & 0.0 & 0.4 & 0.0 & 3.1 & 11.4 & 6.1 \\
$W$, waste flow & Gt & 11.0 & 7.1 & 31.2 & 13.1 & 78.4 & 41.5 & 96.9 & 63.7 \\
$P$, pollution stock & Gt & 79.0 & 39.0 & 368.6 & 188.6 & 1302.0 & 853.3 & 1673.3 & 1224.9 \\
\bottomrule
\end{tabular}
\begin{tablenotes}[flushleft]
\footnotesize
\item[] \textit{Note:} The model column is the planner corner of the calibrated baseline solved from 1900 at $T = 400$; the data column is the model-facing series of the data appendix. Goods are in trillions of 2011 international dollars, material in gigatonnes. The pollution stock is in gigatonnes of material at the bridge composition of the data appendix, which for the data column is the 2000 to 2015 composition at every date, while the model's initial stock is converted at the 1900 composition; the two 1900 entries differ by that conversion alone. The capital stock before 1950 is a perpetual inventory and the material series are the material flow accounts, whose outflow the data column's waste flow is; the model's waste flow also carries unused extraction, which the accounts do not count. Nothing in this table was fitted: every parameter is read from a share, a rate or a single year, and the path is what those parameters imply.
\end{tablenotes}
\end{threeparttable}
\end{table}


\smalltitle{The recycled share.} Table~\ref{tab:q:data:waste} shows the model's recycled share in 2015 more than twice the observed one at a treated share on target, and the two are separable. The handled flow on the solved path in 2015, 66 gigatonnes, is the material flow accounts' whole outflow of that year, processed output plus secondary input, 64 gigatonnes, and the treated share of it matches the treated tonnage on which the observed yield was read. What differs is the yield at the treated tonne: the recycling-capital margin chooses $x=0.43$ and $a=0.69$ where the calibration point of $\xi$ in Table~\ref{tab:q:data:waste} has less than half the capital per treated tonne and a yield of $0.39$, and the factor of $1.8$ between the recovered share of the handled flow on the path and in the accounts is the whole overshoot. The margin over-invests for the reason Section~\ref{sec:q:cal:aggregation} gives: it values a marginal recovered tonne at the mass aggregate's material price, $\gamma Y/R$, between 130 and 180 dollars per tonne in 2015, an average across the fossil, ore and stone prices of Table~\ref{tab:q:data:production}, while the capital outlay through which $\xi$ is set values it at 66 dollars per tonne at an eleven percent return, because the tonnes actually recovered are the cheap ones. The overshoot is the aggregation cost of one material showing up in a number, and not a reason to move $\xi$.

\smalltitle{The metals-only bound.} The bound of Section~\ref{sec:q:data:metals} is infeasible as a path: with $\mu_N=0$ the reserve costate is zero throughout, no exploration takes place, and the reserve goes negative. The metals price channel that put $\mu_N$ there gives a negative elasticity floored at zero, and the floor is arbitrary. The smallest repair with a source is the physical channel already in Table~\ref{tab:q:data:extraction}, the elasticity of energy per tonne to ore grade, set as the metals $\mu_N$ with zero kept as the range's infeasible bottom, and a test that the reserve stays positive before any row is reported; a terminal complementarity on the reserve, which would let a rent exist at $\mu_N=0$, is a change to the model and the second choice. Until one is made the bound's table is not input.

\smalltitle{What state C rests on here.} The classification is by the criteria, not by a resource side that has died. At the horizon extraction is 222 gigatonnes a year, cumulative extraction is $0.58$ of its bound, discovery is large throughout, and the reserve is far from exhausted. The classifier reads the state off the closure criterion and, with a floor, off the survival ratio at the horizon; it does not check that extraction vanishes, and the state-C configuration of \theory{Section}{Perpetual circular growth (C)} has extraction vanishing only in the limit. Two consequences follow for the tables. The retained endowment is read at the horizon and is still accumulating there, so the survival ratio is horizon-dependent: at the largest floor that solves it is an eighth at $T=200$ of its value at $T=400$. Every survival ratio in this section is therefore a lower bound at $T=400$, comparable across cells solved at the same horizon and not a limit. And the level effect of the retained endowment on the balanced path, which \theory{Section}{The circular balanced growth path} names, is in no welfare number here, because every welfare comparison is truncated at the horizon before the resource side dies; at the calibrated preferences it would be discounted away in any case.

\subsection{Where the economy sits in the taxonomy}\label{sec:q:res:taxonomy}

Tables~\ref{tab:q:res:taxonomy} and~\ref{tab:q:res:surface} together answer E1 and E5. The state is decided by the ceiling alone: under the soft ceiling every floor that solves is state C, under either hard ceiling every positive floor is state A and the floorless cell is state B, and there is no interior boundary in the floor anywhere in the solvable range. The surface has no structure in the floor beyond a hyperbola. The retained endowment barely moves along a ceiling row, by two percent between the floorless cell and the largest floor of the grid, and the residence time is constant to four figures, so the survival ratio is a nearly fixed numerator over $\bar R\mathcal T$. What the floor does change is the date the gate fee turns negative, which comes forward as the floor rises, in every ceiling row.

%% GENERATED by model/scripts/run_experiments.jl on 2026-09-18
%% Do not edit by hand: the file is rewritten by the next run.
\begin{table}[!htb]
\centering
\begin{threeparttable}
\caption{The surface over the floor and the ceiling}\label{tab:q:res:surface}
\begin{tabular}{rccrccrccr}
\toprule
$\bar R$ & \multicolumn{3}{c}{$\bar a = 1$} & \multicolumn{3}{c}{$\bar a = 0.85$} & \multicolumn{3}{c}{$\bar a = 0.708$} \\
 & State & $\mathcal M_\infty/(\bar R\mathcal T)$ & $T^{\dagger}$ & State & $\mathcal M_\infty/(\bar R\mathcal T)$ & $T^{\dagger}$ & State & $\mathcal M_\infty/(\bar R\mathcal T)$ & $T^{\dagger}$ \\
\midrule
0.0 & C & $\infty$ & -- & B & $\infty$ & -- & B & $\infty$ & -- \\
2.0 & C & 1001.9 & -- & A & 686.2 & 280 & A & 497.7 & 280 \\
4.0 & C & 500.2 & -- & A & 342.5 & 280 & A & 248.4 & 280 \\
6.0 & C & 333.0 & -- & A & 228.0 & 280 & A & 165.3 & 280 \\
9.502 & C & 209.7 & -- & A & 143.5 & 280 & A & 104.0 & 280 \\
14.0 & C & 141.8 & -- & A & 97.0 & 280 & A & 70.3 & 280 \\
19.0 & C & 104.1 & -- & A & 71.1 & 280 & A & 51.5 & 280 \\
23.756 & C & 83.0 & -- & A & 56.7 & -- & A & 41.0 & -- \\
\addlinespace
edge & C & 56.8 & -- & -- & -- & -- & A & 28.0 & -- \\
\bottomrule
\end{tabular}
\begin{tablenotes}[flushleft]
\footnotesize
\item[] \textit{Note:} Each cell is one solved path at $T = 400$ from the calibrated baseline, the floor and the ceiling overridden. The survival ratio is $\mathcal M_\infty/(\bar R\mathcal T)$ of the theory note's survival condition, infinite at $\bar R = 0$, with $\mathcal M_\infty$ read at the horizon reached and therefore comparable across cells rather than a limit; under a hard ceiling it is reported but does not decide the state, the loop being unable to close. $T^{\dagger}$ is the shutdown date the search returns, and is searched only where the state is collapse. It is not an interior optimum on this calibration: on a ten-period grid the objective rises monotonically in the date and the shutdown branch converges on no date after 280, so the entry is the last date that solves. The last eighty periods of it are worth $3\times 10^{-4}$ of the objective, so the date is better read as `not before 200' than as a date. A dash at $\bar R = 23.756$ is the search converging on no date at the default length of its floor homotopy; at four times the length that cell returns 280 like the others. The floor enters the technology as $R-\bar R$ and 9.5\,Gt is already 1.66 times this calibration's own 1900 material input, so the rows are a grid of different economies rather than a perturbation of the baseline. The last row is the largest floor that solves at all, found by bisecting between the largest floor of the calibration's grid that solves and the smallest that does not, at the two calibrated ceilings only. The edge lies at \(\bar a = 1\) between 34.52 and 34.89\,Gt and at \(\bar a = 0.708\) between 34.52 and 34.89\,Gt, and a finer walk puts it between 34.61 and 34.62\,Gt. Above it no route into the problem converges. The bound is on the base year rather than on the long run: material input takes its minimum in 1900 on every solved path, 2.5 to 2.9\,Gt above the floor, so the edge asks the 1900 economy for about six times the material input it would otherwise choose, and what then fails is the recycling margins' corner conditions on a path that is feasible throughout.
\end{tablenotes}
\end{threeparttable}
\end{table}


\smalltitle{The floor grid is bounded by the base year.} The floors are fractions of 2015 material input, and they are imposed from 1900. On every solved path material input takes its minimum in the base year and grows away from the floor thereafter, so a floor is a demand on the 1900 economy, with the 1900 capital stock, reserve and extraction technology, and on no later date. The largest floor that solves at all, which the note of Table~\ref{tab:q:res:surface} locates, is about six times the model's own 1900 material input, and the plan's fourth grid point, half of 2015 input, is eight times it and out of reach. What fails at the edge is not the floor, which the first period clears, but the recycling margins' corner conditions on a path that is feasible throughout. The surface is therefore a grid of counterfactual 1900 economies, and its edge is a property of the constant floor and of the base year, not of the long run. The survival margin, a ratio of one, lies about fifty times beyond that edge: the floor at which the laissez-faire corner of Section~\ref{sec:q:res:instruments} would fail to clear it is of the order of $1.8\times10^{3}$ gigatonnes, some three hundred times 1900 material input. The version of the question that could be answered needs a floor that scales with the economy, the maintenance reading $\bar R=\bar rK$ that \theory{Appendix}{Extensions and open items} names as the alternative to the constant floor, or, as a smaller change, a constant floor switched on at a date after which the economy is large enough to meet it.

\smalltitle{The collapse cells.} Where a floor meets a hard ceiling the tables report the state and the retained endowment, and in place of a shutdown date they carry the duration bound of \theory{Section}{The material budget and four accounting facts}: cumulative handling cannot exceed the material budget over $1-\bar a$, so at $\bar R$ a period the era lasts at most $\bar{\mathcal R}/\bar R$ periods, which the table notes evaluate. The bound is tens of millennia. Over the four centuries solved, the state-A paths differ from the state-C paths in the ceiling alone: no cell's material era ends within the horizon, and collapse despite growth is a statement about a horizon the solver does not reach. No shutdown date is reported because the search of Section~\ref{sec:q:sol:collapse} returns none: on this calibration its objective is monotone increasing in the date over the whole range on which the branch converges, the increments decay geometrically, the last eighty periods of that range are worth three parts in ten thousand of the objective, and the branch stops converging at about $T^\dagger=285$ for a reason independent of the floor, the handover to the shutdown value function failing with the same residual at every floor. The date the search returns is the horizon of its branch and not a property of the economy, so the tables do not print it. The survival ratio in those cells is the retained endowment over the floor's turnover on a path whose loop cannot close; it decides nothing there, and the tables carry a dash rather than a number that would read as the reason for the state.

\subsection{What circularity buys}\label{sec:q:res:circularity}

Table~\ref{tab:q:res:circularity} reports E2. The no-recycling counterfactual sets the ceiling at $\bar a=0.01$ and keeps everything else, so treatment and the residue leakage of the treated stream remain in the counterfactual economy and the experiment measures recovery, not the waste sector. That the counterfactual's welfare is the same to $2\times10^{-8}$ across the three values of $\xi$ is the check that the tail rate no longer matters at the corner.

%% GENERATED by model/scripts/run_experiments.jl on 2026-09-17
%% Do not edit by hand: the file is rewritten by the next run.
\begin{table}[!htb]
\centering
\begin{threeparttable}
\caption{What circularity buys, by floor, ceiling and tail rate}\label{tab:q:res:circularity}
\begin{tabular}{rrrcrrrrrrrr}
\toprule
$\bar R$ & $\bar a$ & $\xi$ & State & CE gain (\%) & Era ends & $T^{\dagger}$ & $\sum\Xi$ & $\Xi$ avoided & $\mathcal M_\infty$ & $\mathcal M_\infty/(\bar R\mathcal T)$ & Hotelling dev. \\
\midrule
0.0 & 1.0 & 0.91 & C & 0.0009 & -- & -- & 11663.9 & 16767.3 & 74839.2 & $\infty$ & 0.0702 \\
0.0 & 1.0 & 2.72 & C & 0.0065 & -- & -- & 6107.1 & 22311.5 & 86010.3 & $\infty$ & 0.0708 \\
0.0 & 1.0 & 8.17 & C & 0.04 & -- & -- & 3263.4 & 25146.3 & 92343.3 & $\infty$ & 0.0708 \\
0.0 & 0.85 & 0.91 & B & 0.0006 & -- & -- & 16284.7 & 12146.5 & 53246.6 & $\infty$ & 0.0685 \\
0.0 & 0.85 & 2.72 & B & 0.0037 & -- & -- & 11427.4 & 16991.2 & 58780.8 & $\infty$ & 0.0697 \\
0.0 & 0.85 & 8.17 & B & 0.0194 & -- & -- & 8920.1 & 19489.7 & 61641.1 & $\infty$ & 0.0705 \\
0.0 & 0.708 & 0.91 & B & 0.0003 & -- & -- & 19728.4 & 8702.8 & 39710.4 & $\infty$ & 0.0676 \\
0.0 & 0.708 & 2.72 & B & 0.002 & -- & -- & 15180.7 & 13238.0 & 42586.9 & $\infty$ & 0.0683 \\
0.0 & 0.708 & 8.17 & B & 0.0092 & -- & -- & 12689.8 & 15720.0 & 43960.7 & $\infty$ & 0.0686 \\
9.5 & 1.0 & 0.91 & C & 0.001 & -- & -- & 12726.7 & 17266.9 & 74147.5 & 181.9 & 0.0755 \\
9.5 & 1.0 & 2.72 & C & 0.0086 & -- & -- & 6896.3 & 23084.8 & 85387.4 & 209.7 & 0.0755 \\
9.5 & 1.0 & 8.17 & C & 0.095 & -- & -- & 3741.0 & 26231.1 & 91933.6 & 225.8 & 0.0768 \\
9.5 & 0.85 & 0.91 & A & 0.0006 & -- & 280 & 17383.7 & 12609.9 & 52725.0 & 129.6 & 0.0755 \\
9.5 & 0.85 & 2.72 & A & 0.0045 & -- & 280 & 12266.9 & 17714.1 & 58300.5 & 143.5 & 0.0755 \\
9.5 & 0.85 & 8.17 & A & 0.033 & -- & 280 & 9502.3 & 20469.8 & 61263.1 & 150.8 & 0.0755 \\
9.5 & 0.708 & 0.91 & A & 0.0004 & -- & 280 & 20878.5 & 9115.1 & 39306.4 & 96.8 & 0.0755 \\
9.5 & 0.708 & 2.72 & A & 0.0024 & -- & 280 & 16085.2 & 13895.8 & 42207.8 & 104.0 & 0.0755 \\
9.5 & 0.708 & 8.17 & A & 0.0134 & -- & 280 & 13366.4 & 16605.7 & 43638.2 & 107.6 & 0.0755 \\
\bottomrule
\end{tabular}
\begin{tablenotes}[flushleft]
\footnotesize
\item[] \textit{Note:} The no-recycling counterfactual sets $\bar a = 0.01$, so the recycling margin is at its corner rather than absent from the problem, and is solved at the same $\xi$ and the same floor. The consumption-equivalent gain is the proportional consumption supplement that would make the counterfactual as good as the cell. `Era ends' is the first date at which material input reaches the floor and $T^{\dagger}$ the optimal shutdown date, searched only in the collapse cells; a dash is an event that does not occur within the horizon reached. Emissions are cumulative over the horizon, in gigatonnes of material. The survival ratio is infinite without a floor, and under a hard ceiling it is reported but is not what decides the state: the loop cannot close, so a floor puts the cell in state A whatever the ratio. The Hotelling deviation is the largest relative gap between $\Psi_{t+1}/\Psi_t$ and $1+r_{t+1}$ along the path.
\end{tablenotes}
\end{threeparttable}
\end{table}


The leakage effect and the welfare effect are three to four orders of magnitude apart. Across the grid recycling avoids between three tenths and nine tenths of the counterfactual's cumulative leakage, cumulative leakage falls by a factor of six from the worst cell to the best, and the retained endowment more than doubles; the consumption-equivalent gain is below a tenth of a percent in every cell, and the signs are the theory's, rising in $\xi$ and in $\bar a$ and larger with a floor than without. A gain of a hundredth of a percent is within the horizon tolerance of the welfare tail, so the ranking across cells is safer than any single level.

The welfare numbers are a statement about the discount rate and the damage coefficient, and Table~\ref{tab:q:res:sensitivity} says by how much. Three things make them small, and all three are the calibration of Section~\ref{sec:q:data} rather than the solver or the model. The weight the objective puts on a period decays at $\beta e^{(1-\eta)g}$, which at the calibrated preferences is the transversality factor of Section~\ref{sec:q:res:baseline}, so a date two centuries out carries a weight of $8\times10^{-5}$; the pure rate of time preference behind it is the intercept of the seven-point regression of Table~\ref{tab:q:data:preferences} on a gross return to capital. The damage coefficient is a marginal integrated-assessment loss discounted at the seven and a half percent real rate of that table: the price scale $\kappa Y/(r+\theta)$ it puts on a tonne of pollution in 2015 is about 8 dollars, 11 dollars per tonne of CO$_2$ at the bridge of Table~\ref{tab:q:data:damages}, several times below DICE-2023's own social cost of carbon at its own preferences, and the gap is the discount rate again and not the damage function. And the utility channel is off, so damages reach welfare through output alone. The sensitivity table takes the damage function of \textcite{howard2017} with the high end of the transient response, which together multiply the coefficient by the factor its note derives, and the preference pair of \textcite{barrage2024}, separately and together. The gain rises in every column; at the point cell, the only cell the sensitivity runs, it is four times the calibrated value under the larger damage coefficient, thirty times under the DICE-2023 preference pair, and eighty times under both, where it reaches half a percent of consumption. That is the whole range the two least defensible parameters of the calibration span: under no setting in the table is recycling worth more than a percent of consumption over the horizon, and under every setting it avoids most of the counterfactual's leakage. The preference pair is not admissible everywhere in the model: at $\eta=1.45$ the collapse ranking condition of \theory{Section}{Scope: maintained long-run assumptions} needs $\rho$ above $0.0174$ at the calibrated depreciation rate, and DICE-2023's $0.015$ is below it, so a shutdown path has no finite value there and a collapse cell could not be ranked under those preferences. The paths in the table are all state C and do not use the condition; the other four long-run conditions hold at both pairs.

%% GENERATED by model/scripts/run_experiments.jl on 2026-09-18
%% Do not edit by hand: the file is rewritten by the next run.
\begin{table}[!htb]
\centering
\begin{threeparttable}
\caption{The welfare numbers under the alternative damage and preference settings}\label{tab:q:res:sensitivity}
\begin{tabular}{lrrrr}
\toprule
Quantity & Calibrated & Damages $\times6.69$ & DICE-2023 preferences & Both \\
\midrule
\multicolumn{5}{l}{\textit{The setting}} \\
$\kappa$, per Gt & $6.78\times10^{-6}$ & $4.54\times10^{-5}$ & $6.78\times10^{-6}$ & $4.54\times10^{-5}$ \\
$\rho$ & 0.0374 & 0.0374 & 0.015 & 0.015 \\
$\eta$ & 1.409 & 1.409 & 1.45 & 1.45 \\
$\beta e^{(1-\eta)g}$, weight per period & 0.954 & 0.954 & 0.9741 & 0.9741 \\
well-posedness condition failing & none & none & (iii) cake-eating & (iii) cake-eating \\
\addlinespace
\multicolumn{5}{l}{\textit{The planner corner, $\bar a = 1$, $\bar R = 0$}} \\
State & C & C & C & C \\
$\mathcal M_\infty$, Gt & 86010 & 89615 & 93111 & 96853 \\
$\sum\Xi$, Gt & 6107 & 3389 & 5019 & 2346 \\
Gate fee $<0$ from & 139 & 160 & 133 & 151 \\
$\tau^{\mathcal W}_T$ & -5.1 & -4.88 & -9.65 & -9.29 \\
$\mathcal M_\infty/(\bar R\mathcal T)$ at $\bar R = 23.756$ & 83.0 & 86.2 & 86.6 & 90.0 \\
\addlinespace
\multicolumn{5}{l}{\textit{Consumption equivalents, percent}} \\
gain of recycling over $\bar a = 0.01$ & $6.5\times10^{-3}$ & $2.62\times10^{-2}$ & $2.05\times10^{-1}$ & $5.35\times10^{-1}$ \\
cost of laissez-faire $(0,0,0,0)$ & $8.24\times10^{-4}$ & $3.14\times10^{-2}$ & $1.7\times10^{-2}$ & $4.14\times10^{-1}$ \\
cost of the missing emission tax $(1,1,0,1)$ & $7.82\times10^{-4}$ & $3.13\times10^{-2}$ & $1.37\times10^{-2}$ & $4.07\times10^{-1}$ \\
$\sum\Xi$ under laissez-faire, Gt & 7569 & 7684 & 7063 & 7147 \\
$T$ of the laissez-faire comparison & 400 & 358 & 400 & 337 \\
\bottomrule
\end{tabular}
\begin{tablenotes}[flushleft]
\footnotesize
\item[] \textit{Note:} The calibrated column is the baseline of the data appendix. The damage column multiplies $\kappa$ by 6.69, the factor by which the damage function of Howard and Sterner (2017), $\pi_2 = 0.01$ as DICE-2023 reports it, and the high end of the AR6 transient response, 0.00063 against 0.00045 degrees per GtCO$_2$, together raise the marginal loss per gigatonne at the reference stock of the data appendix (3.09 from the damage function alone, 2.02 from the response alone). The preference column is the DICE-2023 pair of Barrage and Nordhaus (2024), $\rho = 0.015$ and $\eta = 1.45$; at that pair the collapse ranking condition $(1-\delta)^{1-\eta} < 1+\rho$ fails at the calibrated depreciation rate, so a shutdown path has no finite value there and collapse cells could not be ranked; the paths on this table are state C and do not use it. The weight per period is the asymptotic factor of the transversality report at the circular growth rate. Every path is the planner corner of the baseline continued in the parameter, the dials or the ceiling, at $T = 400$; $\mathcal M_\infty$ is read at the horizon reached. The consumption equivalents are the proportional consumption supplements of the E2 and E3 tables, measured within each column against that column's own planner corner, which is why a column and not a row is the comparison the table is built for. The laissez-faire corner does not solve to the requested horizon at the larger damage coefficient: continuation in $\kappa$ from the calibrated corner stalls in the treated-share complementarity rows between six and six and a half times the calibrated value, and the cold route converges at the shorter horizon the last row gives; in those columns the corner and the planner corner are compared at that horizon, the planner corner re-solved there, and the cost is read at a horizon at which the welfare tail is a little larger. The gate fee is in trillions of dollars per gigatonne, which is thousands of dollars per tonne.
\end{tablenotes}
\end{threeparttable}
\end{table}


No Hotelling column is reported. The rule of \theory{Section}{The shutdown and its aftermath}, the material value rising at the social discount rate along a collapse path, is a statement about the relaxed problem: the multiplier on the cumulative throughput bound with the period cap dropped, while the floor is slack. On the full problem the material value $\Psi$ carries the reserve costate, which pays the stock-effect dividend $-C^N_S$ whenever $\mu_N>0$ and, while discovery is interior, is pinned by the discovery margin rather than by the interest rate. A comparison of $\Psi_{t+1}/\Psi_t$ with $1+r_{t+1}$ on the full problem must therefore fail, and it fails by seven percent in every cell of the grid, collapse or not, with discovery still at 131 gigatonnes a year in 2100 on the baseline. The check is reported only where the reserve carries no stock effect and exploration has ceased, and no cell reaches that.

\subsection{The cost of the three market failures}\label{sec:q:res:instruments}

Table~\ref{tab:q:res:instruments} reports E3 in two panels, because the question the experiment exists to ask cannot be posed without a floor: the survival ratio is then infinite by construction and no missing instrument can cross it. Panel A measures the cost of the three failures in consumption on the baseline; Panel B, at the largest floor of the grid that solves, is where a crossing could appear.

%% GENERATED by model/scripts/run_d3.jl on 2026-09-17
%% Do not edit by hand: the file is rewritten by the next run.
\begin{table}[!htb]
\centering
\begin{threeparttable}
\caption{The cost of the three market failures, at the corners of the policy space}\label{tab:q:res:instruments}
\begin{tabular}{lcrrrrr}
\toprule
$(\phi^W,\phi^z,\phi^P,\phi^X)$ & State & CE cost ($10^{-5}$\%) & $\mathcal M_\infty/(\bar R\mathcal T)$ & $\mathcal M_\infty$ & $\sum\Xi$ & Gate fee $<0$ \\
\midrule
\multicolumn{7}{l}{\textit{Panel A. $\bar R = 0$, no floor}} \\
(1,1,1,1) & C & 0.0 & $\infty$ & 86010.3 & 13449.8 & 139 \\
(0,1,1,1) & C & 36.66 & $\infty$ & 76332.3 & 12910.3 & -- \\
(1,0,1,1) & C & 7.56 & $\infty$ & 81286.9 & 13526.1 & 139 \\
(1,1,0,1) & C & 78.17 & $\infty$ & 84497.3 & 15214.8 & 0 \\
(1,1,1,0) & C & 0.8 & $\infty$ & 88619.0 & 13888.7 & 139 \\
(0,0,1,1) & C & 36.66 & $\infty$ & 76332.3 & 12910.3 & -- \\
(0,1,0,1) & C & 77.94 & $\infty$ & 74976.2 & 14551.0 & -- \\
(0,1,1,0) & C & 38.44 & $\infty$ & 79100.4 & 13265.4 & -- \\
(1,0,0,1) & C & 81.53 & $\infty$ & 79619.3 & 15507.5 & 0 \\
(1,0,1,0) & C & 8.49 & $\infty$ & 84010.7 & 13932.8 & 140 \\
(1,1,0,0) & C & 84.75 & $\infty$ & 86983.7 & 15766.7 & 0 \\
(0,0,0,1) & C & 77.94 & $\infty$ & 74976.2 & 14551.0 & -- \\
(0,0,1,0) & C & 38.44 & $\infty$ & 79100.4 & 13265.4 & -- \\
(0,1,0,0) & C & 82.38 & $\infty$ & 77637.8 & 15004.9 & -- \\
(1,0,0,0) & C & 87.17 & $\infty$ & 82208.7 & 16038.2 & 0 \\
(0,0,0,0) & C & 82.38 & $\infty$ & 77637.8 & 15004.9 & -- \\
\addlinespace
\multicolumn{7}{l}{\textit{Panel B. $\bar R = 23.756$ Gt, the largest floor that solves}} \\
(1,1,1,1) & C & 0.0 & 82.95 & 84466.9 & 15390.1 & 117 \\
(0,1,1,1) & C & 14.06 & 73.79 & 74971.8 & 14766.7 & -- \\
(1,0,1,1) & C & 8.65 & 78.6 & 79833.7 & 15435.6 & 118 \\
(1,1,0,1) & C & 69.74 & 81.68 & 83168.2 & 16912.0 & 0 \\
(1,1,1,0) & C & 4.34 & 85.5 & 87048.6 & 15815.8 & 118 \\
(0,0,1,1) & C & 14.06 & 73.79 & 74971.8 & 14766.7 & -- \\
(0,1,0,1) & C & 67.0 & 72.62 & 73793.5 & 16199.3 & -- \\
(0,1,1,0) & C & 20.28 & 76.48 & 77700.9 & 15111.6 & -- \\
(1,0,0,1) & C & 73.0 & 77.17 & 78388.0 & 17161.6 & 0 \\
(1,0,1,0) & C & 13.58 & 81.26 & 82521.6 & 15831.1 & 119 \\
(1,1,0,0) & C & 82.95 & 84.11 & 85643.0 & 17436.6 & 0 \\
(0,0,0,1) & C & 67.0 & 72.62 & 73793.5 & 16199.3 & -- \\
(0,0,1,0) & C & 20.28 & 76.48 & 77700.9 & 15111.6 & -- \\
(0,1,0,0) & C & 77.72 & 75.23 & 76428.5 & 16631.7 & -- \\
(1,0,0,0) & C & 85.34 & 79.71 & 80957.1 & 17667.0 & 0 \\
(0,0,0,0) & C & 77.72 & 75.23 & 76428.5 & 16631.7 & -- \\
\bottomrule
\end{tabular}
\begin{tablenotes}[flushleft]
\footnotesize
\item[] \textit{Note:} Each dial is at one, the Pigouvian level of the instrument it switches on, or at zero: $\phi^W$ property rights over the waste stock, $\phi^z$ the material-content charge, $\phi^P$ the emission tax, $\phi^X$ the discovery tax. The consumption-equivalent cost is the proportional consumption supplement that would make the corner as good as the planner corner $(1,1,1,1)$, which is therefore zero by construction, in units of $10^{-5}$ percent of consumption; the planner corner is the maximum of the problem, so a negative entry would be a horizon diagnostic and not a result, and entries below $0.1$ are at the precision of the comparison itself. $\mathcal M_\infty$ is the retained endowment at the horizon reached and $\sum\Xi$ cumulative leakage to the environment, both in gigatonnes; the survival ratio is infinite without a floor. Dates are periods since 1900; a dash is an event that does not occur along the path, and at $\phi^W = 0$ there is no gate fee to change sign, $\tau^{\mathcal W}=-\phi^Wp^{\mathcal W}$ being identically zero. Every row is one solved path at $T = 400$; the four single-dial paths between the corners are in the run's CSVs.
\end{tablenotes}
\end{threeparttable}
\end{table}


\smalltitle{No corner crosses.} Every corner at either floor is state C, so there is no headline corner. At the floor the worst corner clears the survival margin by the factor its survival ratio in Panel B shows, and the three failures move the retained endowment by at most thirteen percent in either direction. Because the endpoint alone could hide a transition, the laissez-faire path is classified period by period against the planner's, the classifier applied to the path truncated at each date: at the floor the two classifications agree at every one of the 401 dates, both read as state A in 1900, where the retained endowment is far below the floor's turnover requirement, and both cross the survival margin in 1921. The laissez-faire ratio is marginally above the planner's until the middle of this century and about nine percent below it thereafter: the leakage the missing instruments cause is a slow stock effect. The floor at which laissez-faire would fail to clear the margin is the one Section~\ref{sec:q:res:taxonomy} gives, and the floors between the solvable edge and it are not reachable with the present solver, so the statement is that the margin is far away on this calibration, not that it has been searched to.

\smalltitle{The discovery tax.} Switching the discovery tax off raises the retained endowment and the survival ratio while lowering welfare, monotonically along the single-dial path, in both panels. That is the accounting of \theory{Section}{Perpetual circular growth (C)}, not a solver artefact: the retained endowment is the initial stocks plus cumulative extraction and displaced mass less cumulative leakage, more exploration raises the inflow term directly, and the mass exploration displaces enters the waste stock, which is retained. The welfare cost of excessive exploration is real resources burnt; its effect on the retained endowment is positive at any finite horizon. The theory note's statement that each failure pushes the margin the wrong way was corrected on this evidence, and \theory{Section}{Laissez-faire and the survival margin} now leaves the direction of the third failure to the quantitative model. One consequence is that laissez-faire is not the worst corner for survival; the corner with the discovery tax alone switched on is.

\smalltitle{Twelve allocations in sixteen corners.} The four pairs of corners that differ only in $\phi^z$ at $\phi^W=0$ coincide to every printed digit in both panels. That is the precondition of \theory{Section}{The three market failures} made exact: with no gate fee the material liability $p^M$ pays no dividend, so $\zeta^{\ast}=\sigma(\tau^W-p^M)$ vanishes and with it $z=\phi^z\zeta^{\ast}$, whatever $\phi^z$. The material-content instrument is inert without property rights over the waste stock, while the emission tax and the discovery tax keep their bases and move the allocation at $\phi^W=0$, as the panels show.

\smalltitle{The ordering check.} The planner corner is the maximum of the problem, so a corner above it would be a horizon diagnostic and not a result. No non-planner point at either floor has a positive consumption-equivalent gain at $T=400$; the entries below $10^{-6}$ percent are at the precision of the comparison, the welfare differences behind them being $10^{-8}$ against welfare levels near forty in magnitude. The same run at $T=40$ puts two corners above the planner, which is the horizon warning of Section~\ref{sec:q:sol:horizon} in the policy space: the equivalents are read at the long horizon only. They are also lower bounds on the cost of a corner: the welfare tail continues consumption at the terminal growth factor and nothing else, and the laissez-faire path ends with slightly more consumption than the planner's, ten percent less retained material and a larger pollution stock, neither of which the tail prices.

\smalltitle{The cost in consumption.} On the calibrated preferences and damages the whole of laissez-faire costs under a hundred-thousandth of consumption, while it raises cumulative leakage by twelve percent and cuts the retained endowment by ten. The missing emission tax is the most expensive single failure at both floors and accounts for most of the laissez-faire cost; with it off the gate fee is negative from the base year, the value Panel A and Panel B print as a zero date, and the single-dial path in $\phi^P$ brings the date forward continuously to it. Under the larger damage coefficient the cost of the missing emission tax rises forty-fold, and under the preference pair the cost of laissez-faire rises twenty-fold; under both the missing emission tax alone costs four tenths of a percent of consumption. Laissez-faire costs more than the missing emission tax in every column, by a margin that is widest under the DICE-2023 pair, where the lower discount rate prices the material the other two failures leak. The laissez-faire corner does not solve to the full horizon at the larger damage coefficient, and its two entries are read at the horizon the table's last row gives, for the reason the note states; the state, the survival ratio at the floor and the ordering against the planner corner are unchanged in every column.

\subsection{Property rights over the waste stock}\label{sec:q:res:gatefee}

Table~\ref{tab:q:res:gatefee} reports E4 at the planner corner and at the corner where property rights are the only instrument. The unit is trillions of dollars per gigatonne, which is thousands of dollars per tonne. At the planner corner the fee is 15 cents per tonne in 1900, peaks at just under two dollars in 2007, crosses zero on the date the table gives, passes one dollar the other way in 2045, ten in 2072, a hundred in 2143 and a thousand in 2237, and ends at the table's terminal value: waste becomes a resource within this century on the point calibration, and a valuable one within the next. Two readings guard the number. The fee is the shadow price of the waste stock and excludes the collection and treatment charges of Table~\ref{tab:q:data:waste}, which is where a real gate fee's money goes; that it is two orders of magnitude below observed tipping fees in the positive phase is what should happen and not a failed validation. And the positive phase is the pollution liability: at the corner where the emission tax is off the fee is negative throughout and numerically zero at the base year, first passing a thousandth of a dollar per tonne in 1958 and a dollar in 2021, and it converges towards the planner's path from below. The sign change the table reports is therefore the date at which the material value of the stock overtakes its pollution liability, and it is the least robust date in this section. Across the range of $\xi$ in Table~\ref{tab:q:data:waste} it moves from 1979 at the high end to 2093 at the low end, 114 years across the one parameter the data block records as its thin spot, monotone and correctly signed, since a cheaper recovery technology makes waste an asset sooner; the terminal fee moves from about 4,900 to 5,500 dollars per tonne over the same range. Across the range of $\mu_N$ the date does not move between the point and the high end. Across the four settings of Table~\ref{tab:q:res:sensitivity} the date moves within three decades. The larger damage coefficient delays it by two decades, because the positive phase is the pollution liability and a higher price on leakage lengthens it; the DICE-2023 preference pair brings it forward by six years only, less than the review of the calibration expected of a discount rate that prices the stock's option value directly, because the same lower rate raises the pollution liability as well; the terminal fee under that pair is nearly twice the calibrated one. The signs along the path are those of \theory{Section}{What the instruments look like along the paths}: the fee positive and $z$ a rebate while the stock is a liability, and after the sign change $p^{\mathcal W}<0$, $\zeta<0$, $\tau^{\mathcal W}<0$ and $z>0$ together, with $\zeta$ turning negative ten periods after the fee.

%% GENERATED by model/scripts/run_experiments.jl on 2026-09-17
%% Do not edit by hand: the file is rewritten by the next run.
\begin{table}[!htb]
\centering
\begin{threeparttable}
\caption{The gate fee and the date waste becomes a resource}\label{tab:q:res:gatefee}
\begin{tabular}{llrrrr}
\toprule
Setting & $(\phi^W,\phi^z,\phi^P,\phi^X)$ & $\tau^{\mathcal W}_0$ & $\tau^{\mathcal W}_T$ & Sign change & $T$ reached \\
\midrule
planner corner & (1,1,1,1) & 0.0002 & -5.0988 & 139 & 400 \\
property rights only & (1,0,0,0) & -0.0 & -3.8487 & 0 & 400 \\
\bottomrule
\end{tabular}
\begin{tablenotes}[flushleft]
\footnotesize
\item[] \textit{Note:} The gate fee is $\tau^{\mathcal W} = -\phi^W p^{\mathcal W}$, a price and not a tax rate: at $\phi^W = 1$ it is the market in which the waste stock is priced at all. The sign change is the first date at which it turns negative, so that holding waste is worth paying for. The per-period paths are in the run's CSV.
\end{tablenotes}
\end{threeparttable}
\end{table}


\subsection{Limitations and open items}\label{sec:q:res:limits}

\smalltitle{The waste stock and its handling share.} The handling share of Table~\ref{tab:q:data:waste} is the disposal flow over its own cumulation, and for a flow growing at a constant rate that ratio is the growth rate: the disposal flow grew at three percent a year over the century, and $\mu$ is its inverse, a residence time that any cumulated series with the same growth would give. It is not identified by the stock-to-flow ratio, and identifying it would need a measured removal from the stock, which no source reports. On the solved path the handled flow in 2015 is the accounts' whole outflow, not the disposal flow the appendix set $\mu$ to reproduce, because the model's stock books every outflow and the unused extraction where the accounts' stock carries disposal alone, and the two readings are consistent only because the stock is larger in the same proportion. The consequence is timing: the accounts emit CO$_2$ and dissipative losses in the year of use, while the model parks every outflow in the stock for a mean residence of three decades, so leakage lags processed output early in the century and its shape differs throughout. What would change it is a stock that distinguishes the flows that are stockpiled from those that are released on use, which is a second waste stock in the model, or a handling share that varies with the composition of the flow.

\smalltitle{Unused extraction in the pollution stock.} The model's transition adds every tonne leaving the waste stock untreated to the pollution stock, and 45 percent of the waste flow in 2015 is unused extraction, which the CO$_2$ bridge of Table~\ref{tab:q:data:damages} never priced: the bridge is the fossil share of processed output, which excludes overburden, so the model prices overburden as if it were a quarter fossil carbon. The fossil share of the model's own leakage is nearer $0.14$ than the $0.26$ the bridge uses, and the pollution stock overshoots the CO$_2$ evidence it was bridged from, as the 2015 column of Table~\ref{tab:q:res:fit} shows. What would change it is either a bridge formed on the model's own waste flow, with unused extraction in the denominator, which halves $\kappa$ and is a refit inside the damage script, or a transition that books unused extraction into the waste stock but not into leakage, which is a change to the model. Neither is made here; the first is the open item.

\smalltitle{The end-point of the extraction cost.} The level $c_{N,\infty}$ of Table~\ref{tab:q:data:extraction} is pinned at the 2006 to 2015 mean of the implied series, and that decade contains the commodity price peak: the mean is $1.9$ times the 2015 value, so the model's 2015 extraction cost is thirteen percent of output against the seven percent cost share it was built to reproduce, a miss the fit's error in logs reports but its end point does not. The range, half of the point, brackets the answer. What would change it is pinning the end on a smoothed price index, a centred ten-year mean of the aggregate, or at the 2015 value, which is a refit inside the extraction script and the open item.

\smalltitle{Four weakly identified parameters.} Four parameters rest on evidence that does not reach the object they parameterise, and each is recorded rather than refitted. The unused-extraction ratio $\Omega^{N,S}$ takes its fossil component, six tenths of the point, from an EU-15 ratio for coal and lignite and applies it to world fossil extraction that is two thirds oil and gas by mass, whose hidden flows are small; it is an upper reading, it is the largest component of the model's waste flow, and the low end of the range in Table~\ref{tab:q:data:accounting} is the better point until the world figure, which is in the literature and not yet obtained, replaces it. The durables coefficient $\phi^I$ is constant against a stored share that tripled over the century, so the model stores too much early and too little late, $0.26$ of its input in 1900 and $0.32$ in 2015 against the data's shares in Table~\ref{tab:q:data:accounting}, and the 1900 waste flow is a tenth low; the consequence is in the residence time and in the split of the retained endowment, not in the state, and a time-varying coefficient is the model change the theory note's setup sets aside. The exploration elasticity $\mu_D$ counts deposits rather than tonnes and is applied to a room below the ceiling that is a fossil-dominated aggregate; it hardly matters at the point, because the ceiling is more than six times the initial cumulative discovery and the exploration cost stays near its floor for centuries, so discovery is cheap and large; the sensitivity that would matter is the ceiling at its low end, $1.27$ times the initial value, which no run touched and which is the open item. And $\rho$ is the intercept of a seven-point regression on a gross return to capital, which Section~\ref{sec:q:res:circularity} shows to decide every welfare level; the sensitivity table is what stands in for its identification.

\smalltitle{The horizon dependence of the retained endowment.} The retained endowment is read at $T=400$ on paths whose resource side is alive, so every survival ratio is a lower bound at that horizon and the comparisons across cells and corners are at the same horizon by construction. On a state-C path the endowment keeps accumulating until extraction vanishes, and the classification does not test that it does. What would change it is a longer horizon, which the conditioning wall of Section~\ref{sec:q:sol:horizon} permits on this calibration since the reserve is far from exhausted, or a terminal closure that imposes the state-C limit on the retained endowment and reports the ratio there.

\smalltitle{The composition of the initial pollution stock.} $P_0$ is the 1900 atmospheric excess stock converted at the 1900 composition of the outflow, while the model-facing pollution series is converted at the 2000 to 2015 composition and $\kappa$ is bridged at the same later composition, so the 1900 entries of Table~\ref{tab:q:res:fit} differ by a factor of two on the conversion alone and $\kappa P_0$ prices the 1900 stock as twice the CO$_2$ it contains. The loss at the initial stock is small either way, but the two objects should share one composition. What would change it is converting $P_0$ at the composition $\kappa$ is bridged at, which halves it and is a one-line change in the pipeline, with the appendix saying that the pollution stock is in tonnes at one stated composition.
